\section{Continuity of the nonlinear output on a fixed torus}
\label{sec:finite-volume-parameter-continuity}
The nonlinear step must vary continuously with its incoming coupling
and activity array. This assertion concerns the actual array in a
fixed support dictionary. Equality of the evaluated gas would be
insufficient. The constants in this section may depend on the fixed
torus; uniform small prices for the scale construction are established
separately.

\subsection{Which data are held fixed}
Consider two incoming representations $R,R'$ on the same fine torus.
They have the same regular and forced addresses, thresholds, clipping
convention and path data. Write $\beta,\beta'>0$ for their couplings.
Their regular and forced activity differences have nonnegative
amplitudes $a_i,b_i$ with
\[
                  \sum_i a_i\le t,\qquad \sum_i b_i\le t.
\]
The regular differences here are differences of the actual masked
occurrences. Their supports and healthy masks agree exactly.
Let $g(\beta)=\beta^{-1/2}$ and put
$\lambda=g(\beta)/g(\beta')>0$. The remaining comparison conditions
are
\begin{equation}\label{eq:parameter-close-input}
 |\lambda^2-1|\le t,\quad |\beta'-\beta|\le t,\quad
 |\log m(\beta')-\log m(\beta)|\le t,
 \qquad 0<t\le\min(1/4,\beta/2),
\end{equation}
where $m(\beta)>0$ is the fixed local Haar-chart normalizing factor.
All forests, tags, component labels and output address identifications
are retained. In particular, their compatibility and eligibility
conditions agree under the identity on elementary labels.

For a continuous incoming curve these conditions hold locally for
arbitrarily small $t$. The activity conditions follow from its array
continuity; multiplication by the fixed healthy mask preserves the
value difference and makes both occurrences zero off that mask.
The scalar conditions follow from continuity of the positive coupling,
its inverse square root, and the logarithm of $m$. This observation
uses the same path throughout the curve.

\subsection{A Gaussian rescaling estimate}
We explain why changing the coupling does not require differentiating
the incoming measurable activities. On one fixed Gaussian space of
dimension $d$, with positive precision matrix $Q$, put
\[
 p(z)=c_Q e^{-z^TQz/2},\qquad
 J_\lambda(z)=\lambda^d\frac{p(\lambda z)}{p(z)}.
\]
Thus, if $C$ denotes the Gaussian covariance, $Q=C^{-1}$; the two
matrices are kept distinct throughout this comparison.
Then for every integrable function $F$,
\begin{equation}\label{eq:parameter-gaussian-rescale}
       \int F(z)\,d\gamma_Q(z)
       =\int J_\lambda(z)F(\lambda z)\,d\gamma_Q(z).
\end{equation}
This is the ordinary change of variables $w=\lambda z$. The Haar
factor, when present, is left unchanged. Since
$g(\beta')\lambda=g(\beta)$, the rescaled incoming field has the
same physical noise; the corresponding compensation map keeps the
original insertion values fixed.

For $|\lambda^2-1|\le t\le1/4$, we have $\lambda\ge1/2$ and
$|\log\lambda|\le t$. Consequently
\[
 |\log J_\lambda(z)|\le
                      t\{d+\tfrac12\|Q\|\,|z|^2\}.
\]
Let $F_0(z)=1+d+\sum_{r=1}^d|z_r|$ and choose a finite constant
$C_Q$ such that the braces are at most $C_QF_0(z)^2$.
The inequality $|u-1|\le(1+u)|\log u|$ for $u>0$ follows by
integrating $e^s$ between $0$ and $\log u$. Thus
\[
 |J_\lambda-1|\le t C_Q(1+J_\lambda)F_0^2.
\]
All polynomial moments of $F_0$ are finite, by the finite Gaussian
moment argument. If $M_k=\int F_0^k\,d\gamma_Q$, then
\eqref{eq:parameter-gaussian-rescale} and
$F_0(\lambda^{-1}z)\le2F_0(z)$ give
\begin{align}
 \int J_\lambda F_0^k\,d\gamma_Q&\le2^kM_k,
 \int |J_\lambda-1|F_0^k\,d\gamma_Q
                   &\le t C_Q(1+2^{k+2})M_{k+2}.
 \label{eq:parameter-gaussian-moment-difference}
\end{align}
The same estimates hold for the full finite collection of Gaussian
read coordinates and the product with normalized Haar measure.
No smaller marginal is substituted when the parameter changes.

Here is the precise pointwise input needed to integrate local jet
differences. Fix integers $j,k\ge0$. For every component $P$, suppose its compensated
integrands $f_P,f_P'$ have, on each admissible local variation, the
bound
\begin{equation}\label{eq:parameter-integrand-difference}
 \|J_\lambda f_P'(\cdot,\lambda z)-f_P(\cdot,z)\|_{J^j}
 \le t c_P F_0(z)^k
                 \{J_\lambda(z)+t^{-1}|J_\lambda(z)-1|+1\},
 \qquad \sum_P c_P\le C.
\end{equation}
Here $c_P\ge0$. The constants and $k$ are common to the neighboring representation
and to the forests and tags on this fixed torus. This is the
pointwise compensated comparison estimate, a stronger input than
individual boundedness of the two integrands. Its role is explicit:
\eqref{eq:parameter-gaussian-moment-difference} gives uniform integral bounds. For a pointwise dominator, use
\[
 J_\lambda(z)\le(5/4)^{d/2}\exp(z^TQz/8).
\]
Together with $t^{-1}|J_\lambda-1|\le C_Q(1+J_\lambda)F_0^2$,
this bounds the integrand after division by $t$ by a fixed polynomial
times $1+\exp(z^TQz/8)$. Multiplying by the Gaussian density leaves
a decaying exponential $\exp(-3z^TQz/8)$, so this dominator is
integrable. Integration
and the derivative identities of
Section~\ref{sec:gaussian-differentiation} give
\begin{equation}\label{eq:parameter-component-difference}
 \|F_P'-F_P\|_{J^j}\le t C_1c_P,
                         \qquad \sum_P C_1c_P\le C_1C.
\end{equation}
For example one may take
$C_1=(1+2^k)M_k+C_Q(1+2^{k+2})M_{k+2}$.
The same bound survives normalized forest and tag averaging and
the common local extension, since their germs agree on the healthy
domain. The construction of \eqref{eq:parameter-integrand-difference}
from the individual logarithmic, packet and compensation factors
remains a separate local estimate; it is not inferred merely from
continuity of the scalar coupling.

\subsection{A summable difference through the cluster logarithm}
We prove the summation step rather than interchange an infinite
series with a limit without a common bound. Write $u_P\ge0$ for
component jet prices. Fix a target support exponent $a\ge0$; the cluster comparison itself also holds for arbitrary real $a$. Suppose the ordinary cluster criterion holds
for $2^ju_P$, with reserve $\delta>0$, root price $\gamma>0$ and
ratio $1/2$. Suppose component differences have prices $t d_P$,
where $d_P\ge0$ and $\sum_P d_P\le C$. On the fixed torus let
\[
 K_0=\sum_{Y\subseteq\Lambda}e^{(a+\delta)c(Y)},
 \qquad
 \kappa=\frac{\gamma}{2K_0 2^j(1+C)}>0.
\]
The finite sum includes every possible support. For every root, the
additional vertex budget from $2^j\kappa d_P$ is at most
$K_0 2^j\kappa C\le\gamma/2$. The augmented prices
$2^j(u_P+\kappa d_P)$ therefore satisfy the cluster criterion with
root price $3\gamma/2$ and ratio $3/4$: the degree quantity is linear
in the root price. This chooses one neighborhood size $\kappa$
after the torus and old criterion are fixed.

For $0\le t\le\kappa$, set $\theta=t/\kappa\le1$. The finite product
identity gives, in the joint jet bounds,
\begin{equation}\label{eq:parameter-product-difference}
 \left\|\prod_{r=1}^qF_{P_r}'-\prod_{r=1}^qF_{P_r}\right\|_{J^j}
 \le\theta(1+q)^j
       \left\{\prod_{r=1}^q(u_{P_r}+\kappa d_{P_r})
                                  -\prod_{r=1}^q u_{P_r}\right\}.
\end{equation}
To verify it, expand each $F_P'=F_P+(F_P'-F_P)$. Every nonempty
subset of difference factors contains at least one factor $\theta$;
its higher powers are bounded by $\theta$. At derivative order $r$
there are $q^r$ assignments to the factors, bounded by $(1+q)^j$.
Summing the nonempty subsets gives exactly the difference of the
positive products on the right. Repeated component labels are still
distinct positions in this calculation.

Since $1+q\le2^q$ for $q\ge1$, the derivative factor is paid by
$2^j$ per component. Multiply by the absolute Ursell coefficient
and its factorial normalization, then use the augmented cluster
criterion. Summing rooted bounds over all $|\Lambda|$ sites bounds
the unrooted sum, because every nonzero cluster has nonempty support.
Thus cluster differences have nonnegative prices $t e_\Gamma$ with
\[
 \sum_\Gamma e_\Gamma
       \le\kappa^{-1}|\Lambda|\frac{3\gamma/2}{1-3/4}.
\]
This estimate uses the original cluster criterion with a paid
perturbation margin; no new uniform small-price theorem for the
neighboring representation has been assumed.

\subsection{Forced products and the output array}
The same product calculation applies to the forced exponential
expansion. If its component values have summable bounds $h_i$ and
differences bounded by $t d_i$, with $0\le t\le1$, then a finite
multiplicity vector $\nu$ satisfies
\[
 \left|\prod_i\frac{(F_i')^{\nu_i}}{\nu_i!}
             -\prod_i\frac{F_i^{\nu_i}}{\nu_i!}\right|
       \le t\prod_i\frac{(h_i+d_i)^{\nu_i}}{\nu_i!}.
\]
Indeed apply the same difference-of-products argument to repeated
factors and divide by their positive factorials. Summing over all
finite multiplicity vectors yields
\[
 \sum_\nu\prod_i\frac{(h_i+d_i)^{\nu_i}}{\nu_i!}
                         =\exp\!\left(\sum_i(h_i+d_i)\right),
\]
by finite products of the exponential series followed by monotone
convergence. Compatibility or eligibility conditions only remove
terms; they agree under the fixed address identification. This proves
a summable forced difference bound.

The remaining fixed path terms are affine in the coupling. If
$S_i(\beta)$ is such a term, then
$S_i(\beta)-S_i(\beta')=(\beta-\beta')\{S_i(1)-S_i(0)\}$.
Their two summable derivative prices therefore bound the difference
by $t$ times a summable array. A finite symmetry average preserves
summability after its explicit fixed derivative factor. Any support
weight $e^{Ac(Y)}$ on this fixed torus is at most $e^{A|\Lambda|}$,
and moving an anchor costs at most a fixed power of
$1+\operatorname{diam}\Lambda$. Hence all these operations convert
the preceding sums into the prescribed regular array budget in the jet norm and forced array
budget in the value supremum norm, bounded by $Kt$ for a finite $K$ depending on the base
representation and torus.

This also supplies the common majorant needed along a subsequence.
If input parameters $x_n$ tend to $x$, select an increasing subsequence
$n_k$ for which the output differences from $x$ have summable
nonnegative prices $d_{k,i}$ with $\sum_i d_{k,i}\le2^{-k}$.
If $h_i$ majorizes the output at $x$, put
\[
                         H_i=h_i+\sum_k d_{k,i}.
\]
Tonelli's theorem gives $\sum_iH_i<\infty$, and the triangle
inequality shows that $H_i$ bounds both the output at $x$ and all
outputs on this subsequence, component by component in the required
jet or value norm. The regular and forced arrays are treated
separately in this construction. It proves the required continuity
and subsequential domination in the common dictionary, conditional
on the local compensated difference estimate
\eqref{eq:parameter-integrand-difference} and its forced counterpart.
